% ID: SEL-D0019-OB-E01
% ESTADO: validada. S020, cotejada con examen y criterios oficiales.
\textbf{a)} $A$ es $2\times2$ y $B^t$ es $3\times2$.
Para que los índices interiores coincidan, $P$ ha de ser
$\boxed{2\times3}$, con resultado $2\times2$.
Como $AC$ es $2\times3$, $Q$ ha de tener dos columnas y tres filas
para obtener una matriz $3\times3$. Por tanto $Q$ es $\boxed{3\times2}$.

\textbf{b)} $\det A=2$, así que multiplicamos $AX=A^2+2BC^t$
por $A^{-1}$ a la izquierda. Así,
\[
BC^t=\begin{pmatrix}-2&-5\\11&4\end{pmatrix},\qquad
X=A+\begin{pmatrix}4&-2\\3&-1\end{pmatrix}
\begin{pmatrix}-2&-5\\11&4\end{pmatrix}
=\boxed{\begin{pmatrix}-31&-26\\-20&-15\end{pmatrix}}.
\]
Para comprobarlo, $AX=\begin{pmatrix}-9&-4\\13&18\end{pmatrix}$.
Al restar $2BC^t$ queda $\begin{pmatrix}-5&6\\-9&10\end{pmatrix}=A^2$.
